%=============================================
% Date: August 26, 2026
% Topic: deriving relativity without light
%=============================================

\subsection{Relativity without Light}
	We begin this section by doing a derivation of special relativity 
	in large contrast with standard methods which introduce relatvity
	with light-based thought experiments. Our derivation follows the 
	elegant logic developed by von Ignatowsky (1910) and later refined
	by others, showing that the Lorentz transformations follow from 
	fundamental symmetries alone.  

	\subsubsection{The Transformation Problem}
		Consider two inertial reference frames $S$ and $S'$, where $S$
		is taken as a \textit{rest} frame of reference, and $S'$ moves 
		with a constant velocity $v$ along the $x$ axis with respect to
		$S$. We can express these as follows:
		
		In Galilean transformations, the relationship betwen the two reference
		frames can be described as 
		\begin{align*}
			x' &= x + vt\\
			t' &= t'
		\end{align*}
		That is, expressed in matrix form, 
		\begin{equation}
			\bmqty{
				x'\\t'
			}
			= 
			\bmqty{
				1 & v\\
				0 & 1
			}
			\bmqty{
				x\\t
			}
		\end{equation}
		That implies that time $t$ is absolute, and, the reference frames
		$S$ and $S'$ are not strictly equivalent.
		
		As such, we need a more robust transformation. We seek the most general 
		transformation between the coordinates of the two frames.
		\begin{align}
		    x' = \overline X(x,t;v)\\
			t' = \overline T(x,t;v)
			\label{eq:general_transform}
		\end{align}
		The functions $\overline X$ and $\overline T$ are transformation functions that are initially unknown and undetermined.
		Our task, then, is to determine them from general physical principles. Eventually, our goal
		is to obtain a transformation in the form of Lorentzian transfromations.

		\begin{equation}
			\bmqty{
				x'\\ t'
			}
			=
			\gamma
			\bmqty{
				1 & -v \\
				-\beta & 1
			}
			\bmqty{
				x\\t
			}
		\end{equation}
		where $\gamma = \dfrac{1}{\sqrt{1-\beta^2}}$ and $\beta = \dfrac{v^2}{c^2}$

		An interesting question arises: Where does $c$ come from if we never assume light?

		\subsubsection{Principle of relativity}

		The first assumption we make is the principle of relativity.
		\begin{quote}
		    The laws of physics take the same form in all inertial reference frames.
		\end{quote}

		In particular, $S$ and $S'$ are physically equivalent. There is no reason to regard 
		the transformation from $S$ to $S'$ as fundamentally different from the transfomration 
		from $S'$ back to $S$.

		If $S'$ moves at velocity $v$ relative to $S$, then $S$ moves at a velocity $-v$ relative 
		to $S'$. Hence, the inverse transformation of \ref{eq:general_transform} must have the same functional form:
		\begin{align*}
		    x' = \overline X(x,t;-v)\\
			t' = \overline T(x,t;-v)
		\end{align*}

		This implies that if we apply the forward transformation followed by the inverse transformation, 
		we recover the original coordinates:
		\begin{equation}
			\begin{aligned}
				x &= \bar{X}\bigl(\bar{X}(x,t,v), \bar{T}(x,t,v), -v\bigr) \\
				t &= \bar{T}\bigl(\bar{X}(x,t,v), \bar{T}(x,t,v), -v\bigr)
			\end{aligned}
			\label{eq:composition-identity}
		\end{equation}

		\subsubsection{Isotropy of Space}

		The isotropy of space means that physics has no preferred direction. 
		In particular, the physics cannot distinguish between the $+x$ and $-x$ directions.
		This imposes the following:

		Consider reversing the spatial axis,
		\begin{equation}
			x\rightarrow -x.
		\end{equation}
		Under this reversal, the relative velocity also changes sign,
		\begin{equation}
			v\rightarrow -v.
		\end{equation}

		The spatial coordinate must therefore transform as
		\begin{equation}
			x'=-\bar X(x,t;v),
		\end{equation}

		while time, being a scalar under spatial reflection, satisfies
		\begin{equation}
			t=\bar T(x,t;v).
		\end{equation}
		This distinction is important: $x$ changes sign under spatial reflection, but $t$ does not.

		Under time reversal combined with spatial reflection:

		\begin{equation}
			\begin{aligned}
				x' &= -\bar{X}(-x, t, -v) \\
				t' &= \bar{T}(-x, t, -v)
			\end{aligned}
			\label{eq:isotropy-conditions}
		\end{equation}

		For these to be consistent with the original transformations, we require:

		\begin{equation}
			\begin{aligned}
				\bar{X}(-x, t, -v) &= -\bar{X}(x, t, v) \\
				\bar{T}(-x, t, -v) &= \bar{T}(x, t, v)
			\end{aligned}
			\label{eq:parity-conditions}
		\end{equation}
	\subsubsection{The Invariance of Straigt Lines}

